\chapter{Introduzione}
\label{cap:introduzione}

%Introduzione al contesto applicativo.\\
%
%\noindent Esempio di utilizzo di un termine nel glossario \\
%\gls{api}. \\
%
%\noindent Esempio di citazione in linea \\
%\cite{site:agile-manifesto}. \\
%
%\noindent Esempio di citazione nel pie' di pagina \\
%citazione\footcite{womak:lean-thinking} \\

\section{L'azienda}

RiskApp S.r.l. (Figura \ref{fig:riskapp}) è un'azienda con sede a Conselve (PD) che si occupa di sviluppo software per il mondo assicurativo.\newline
È stata fondata nel 2016 e il suo \textit{core business} è lo sviluppo e il mantenimento dell'omonima applicazione, che viene costantemente aggiornata ed estesa per garantire un prodotto che possa rispondere ad ogni esigenza.\newline
Il principale punto di forza di questa piattaforma è quello di stimare le possibili perdite economiche di un’impresa attraverso un algoritmo proprietario che, anche attraverso l’uso dell'intelligenza artificiale, valuta il rischio raccogliendo e combinando una moltitudine di dati da diverse fonti.\newline
Il personale aziendale lavora costantemente per migliorare i propri servizi, ragionando sui possibili problemi che l'utente e le aziende possono andare incontro, fanno riunioni e call per capire come migliorare e ampliare la piattaforma, tutto svolto in un clima di calma e rispetto tra colleghi.

\begin{figure}[!h] 
    \centering 
    \includegraphics[width=0.4\columnwidth]{RiskAPP} 
    \caption{Logo dell'azienda RiskApp}\label{fig:riskapp}
\end{figure}

\newpage

\section{L'idea}

Per poter gestire le spese per i pasti, che preparano in azienda, è stato scelto di sviluppare un'app mobile che permetta di monitorare i versamenti degli utenti, scegliere il piatto del giorno da un menu condiviso e monitorare la \emph{\gls{cassacomuneg}}\glsfirstoccur.\newline
Deve essere gestita l'autenticazione di ogni utente, dividendo tra utente semplice e utente amministratore e permettere il controllo delle presenze in azienda durante i pranzi. \newline
Ogni utente potrà aggiungere un piatto nel menu, proporre il pasto del giorno, monitorare la sua \emph{\gls{quotastornatag}}\glsfirstoccur e la cassa comune, indicare le spese effettuate e modificare i dati personali. \newline
L'amministratore potrà anche gestire le presenze e le spese effettuate dagli stagisti.
L'applicazione dovrà essere sviluppata con \emph{\gls{flutterg}}\glsfirstoccur, \emph{\gls{dartg}}\glsfirstoccur  e \emph{\gls{firebaseg}}\glsfirstoccur.

\section{Obiettivi}
\label{sec:obbiettivi}

Si farà riferimento ai requisiti secondo le seguenti notazioni:
\begin{itemize}
    \item O per i requisiti obbligatori, vincolanti in quanto obiettivi primari richiesti dal committente;
    \item D per i requisiti desiderabili, non vincolanti o strettamente necessari, ma dal riconoscibile valore aggiunto;
    \item F per i requisiti facoltativi, rappresentanti valore aggiunto non strettamente competitivo.
\end{itemize}
Le sigle precedentemente indicate saranno seguite da una coppia sequenziale di numeri, identificativo del requisito.\newline

\begin{table}[htb]%
    \caption{Tabella degli obiettivi}
    \label{tab:obiettivi-richiesti}
    \begin{tabularx}{\textwidth}{lX}
    \hline
    \textbf{Obiettivo} & \textbf{Descrizione} \\
    \hline\hline
    O01     & Accesso tramite credenziali \\
    \hline
    O02     & Pannello di controllo degli utenti \\
    \hline
    O03     & Modifica o aggiunta delle spese \\
    \hline
    O04     & Monitoraggio della cassa comune \\
    \hline
    O05     & Controllo presenza in azienda di una persona durante i pranzi \\
    \hline
    O06     & Impostare il menu del giorno \\
    \hline
    O07     & Scelta di un piatto dal menu \\
    \hline
    D01     & Integrazione di ChatGPT per consigliare alcune ricette \\
    \hline
    F01     & Test a livello di frontend \\
    \hline
    \end{tabularx}
\end{table}%

\newpage

\section{Organizzazione del testo}

\begin{description}
    \item[{\hyperref[cap:processi-metodologie]{Il secondo capitolo}}] descrive in che modo è stato creato il prodotto desiderato, quale metodo di sviluppo è stato utilizzato e quali sono le tecnologie adottate per lavorare al progetto.
    
    %\item[{\hyperref[cap:descrizione-stage]{Il terzo capitolo}}] approfondisce ...
    
    \item[{\hyperref[cap:analisi-requisiti]{Il terzo capitolo}}] approfondisce i requisiti con una analisi dettagliata di cosa è stato richiesto.
    
    \item[{\hyperref[cap:progettazione-codifica]{Il quarto capitolo}}] approfondisce la progettazione, i \emph{design pattern} utilizzati e la struttura del codice.
    
    %\item[{\hyperref[cap:verifica-validazione]{Il quinta capitolo}}] approfondisce ...
    
    \item[{\hyperref[cap:conclusioni]{Nel quinto capitolo}}] vengono riportate le valutazioni e le conclusioni personali del prodotto.
\end{description}


%da tenere e non toccare --------------------------------------------------------------

Riguardo la stesura del testo, relativamente al documento sono state adottate le seguenti convenzioni tipografiche:
\begin{itemize}
	\item gli acronimi, le abbreviazioni e i termini ambigui o di uso non comune menzionati vengono definiti nel glossario, situato alla fine del presente documento;
	\item per la prima occorrenza dei termini riportati nel glossario viene utilizzata la seguente nomenclatura: \emph{parola}\glsfirstoccur;
	\item i termini in lingua straniera o facenti parti del gergo tecnico sono evidenziati con il carattere \emph{corsivo}.
\end{itemize}
